% =====================================================================
%  minggu-03.tex  --  Slide Pertemuan Minggu ke-3
%  Pengantar Kriptografi (MPFA014) -- Prodi Matematika FMIPA UBB
%
%  Topik silabus : Landasan matematika
%  Sumber materi : buku_ajar/chapters/03 (Bab 3)
%  Sub-CPMK      : 2.3
%
%  Kompilasi dari dalam folder ppt/:
%      compile-ppt.bat 03
% =====================================================================

\documentclass[aspectratio=169,11pt]{beamer}
% =====================================================================
%  beamer-preamble.tex  --  prasarana dan gaya slide kuliah
%
%  Mata kuliah : Pengantar Kriptografi (MPFA014)
%  Program     : Program Studi Matematika, FMIPA, Universitas Bale Bandung
%
%  Berkas ini TIDAK memuat \documentclass. Setiap berkas minggu memuatnya
%  sendiri lalu menarik berkas ini:
%
%      \documentclass[aspectratio=169,11pt]{beamer}
%      % =====================================================================
%  beamer-preamble.tex  --  prasarana dan gaya slide kuliah
%
%  Mata kuliah : Pengantar Kriptografi (MPFA014)
%  Program     : Program Studi Matematika, FMIPA, Universitas Bale Bandung
%
%  Berkas ini TIDAK memuat \documentclass. Setiap berkas minggu memuatnya
%  sendiri lalu menarik berkas ini:
%
%      \documentclass[aspectratio=169,11pt]{beamer}
%      % =====================================================================
%  beamer-preamble.tex  --  prasarana dan gaya slide kuliah
%
%  Mata kuliah : Pengantar Kriptografi (MPFA014)
%  Program     : Program Studi Matematika, FMIPA, Universitas Bale Bandung
%
%  Berkas ini TIDAK memuat \documentclass. Setiap berkas minggu memuatnya
%  sendiri lalu menarik berkas ini:
%
%      \documentclass[aspectratio=169,11pt]{beamer}
%      \input{beamer-preamble}
%      \title{...} \subtitle{...} \setWeekInfo{Minggu n}{Topik}
%
%  Berkas minggu hanya berisi \setWeekInfo dan lingkungan frame.
%  Kompilasi dijalankan DARI DALAM folder ppt/ (lihat compile-ppt.bat),
%  karena gambar buku dirujuk dengan jalur relatif ../buku_ajar/.
% =====================================================================

% ---------------------------------------------------------------- bahasa
\usepackage[utf8]{inputenc}
\usepackage[T1]{fontenc}
\usepackage[indonesian]{babel}
\usepackage{lmodern}
\usepackage{microtype}

% ------------------------------------------------------------- matematika
\usepackage{amsmath,amssymb,mathtools}

% ----------------------------------------------------------------- tabel
\usepackage{booktabs}
\usepackage{array}
\usepackage{tabularx}

% --------------------------------------------------------------- grafik
\usepackage{tikz}
\usetikzlibrary{arrows.meta,positioning,shapes.geometric,calc,fit,backgrounds}

% ---------------------------------------------------------- kode program
\usepackage{listings}
\lstdefinestyle{slidecode}{
    basicstyle=\ttfamily\scriptsize,
    keywordstyle=\color{ubbblue}\bfseries,
    commentstyle=\color{ubbgreen}\itshape,
    stringstyle=\color{ubbred},
    numbers=left,
    numberstyle=\tiny\color{gray},
    numbersep=6pt,
    frame=single,
    rulecolor=\color{ubbblue!40},
    backgroundcolor=\color{ubbblue!3},
    breaklines=true,
    showstringspaces=false,
    tabsize=4,
    columns=fullflexible,
    keepspaces=true,
}
\lstset{style=slidecode}

% =====================================================================
%  WARNA INSTITUSI
% =====================================================================
\definecolor{ubbblue}{RGB}{12, 62, 120}      % warna utama
\definecolor{ubbteal}{RGB}{0, 118, 110}      % aksen kedua
\definecolor{ubbgreen}{RGB}{22, 122, 74}     % untuk blok contoh
\definecolor{ubbred}{RGB}{168, 42, 42}       % untuk blok peringatan
\definecolor{ubbgray}{RGB}{70, 76, 84}       % teks sekunder

% =====================================================================
%  IDENTITAS MATA KULIAH
% =====================================================================
\newcommand{\CourseName}{Pengantar Kriptografi}
\newcommand{\CourseCode}{MPFA014}
\newcommand{\CourseSemester}{Semester Ganjil}
\newcommand{\CourseProgram}{Program Studi Matematika}
\newcommand{\CourseFaculty}{Fakultas Matematika dan Ilmu Pengetahuan Alam}
\newcommand{\CourseInstitution}{Universitas Bale Bandung}
\newcommand{\CourseLecturer}{Cecep Suwanda, S.Si., M.Kom.}

% Label minggu, disetel tiap berkas lewat \setWeekInfo
\newcommand{\WeekLabel}{}
\newcommand{\WeekTopic}{}
\newcommand{\setWeekInfo}[2]{%
    \renewcommand{\WeekLabel}{#1}%
    \renewcommand{\WeekTopic}{#2}%
}

% =====================================================================
%  TEMA: header, footer, warna blok
% =====================================================================
\usetheme{default}
\setbeamertemplate{navigation symbols}{}
\setbeamertemplate{blocks}[rounded][shadow=false]

\setbeamercolor{palette primary}{bg=ubbblue, fg=white}
\setbeamercolor{palette secondary}{bg=ubbgray, fg=white}
\setbeamercolor{palette tertiary}{bg=ubbteal, fg=white}

\setbeamercolor{structure}{fg=ubbblue}
\setbeamercolor{normal text}{fg=black!92}
\setbeamercolor{alerted text}{fg=ubbred}

\setbeamercolor{block title}{bg=ubbblue, fg=white}
\setbeamercolor{block body}{bg=ubbblue!6, fg=black!92}
\setbeamercolor{block title alerted}{bg=ubbred, fg=white}
\setbeamercolor{block body alerted}{bg=ubbred!6, fg=black!92}
\setbeamercolor{block title example}{bg=ubbgreen, fg=white}
\setbeamercolor{block body example}{bg=ubbgreen!6, fg=black!92}

\setbeamercolor{itemize item}{fg=ubbblue}
\setbeamercolor{itemize subitem}{fg=ubbteal}
\setbeamercolor{enumerate item}{fg=ubbblue}

\setbeamerfont{headline}{size=\scriptsize}
\setbeamerfont{footline}{size=\scriptsize}
\setbeamerfont{block title}{size=\normalsize}

% Badan blok rata kiri (bukan rata kanan-kiri): teks slide lebih mudah dibaca
% dan baris terakhir yang pendek tidak lagi memicu "Underfull \hbox".
\addtobeamertemplate{block body begin}{\raggedright}{}
\addtobeamertemplate{block alerted body begin}{\raggedright}{}
\addtobeamertemplate{block example body begin}{\raggedright}{}

% Header: nama mata kuliah dan kode MK
\setbeamertemplate{headline}{%
    \begin{beamercolorbox}[wd=\paperwidth,ht=2.6ex,dp=1.2ex,leftskip=1em,rightskip=1em]{palette primary}%
        \usebeamerfont{headline}\CourseName\ \textperiodcentered\ \CourseCode%
        \hfill \CourseProgram%
    \end{beamercolorbox}%
}

% Footer: minggu dan nomor halaman
\setbeamertemplate{footline}{%
    \begin{beamercolorbox}[wd=\paperwidth,ht=2.4ex,dp=1.2ex,leftskip=1em,rightskip=1em]{palette secondary}%
        \usebeamerfont{footline}\WeekLabel\ \textperiodcentered\ \WeekTopic%
        \hfill \insertframenumber\,/\,\inserttotalframenumber%
    \end{beamercolorbox}%
}

% Halaman judul: judul, subjudul, lalu blok identitas mata kuliah
\setbeamertemplate{title page}{%
    \begin{centering}
        \vskip1.0em
        {\usebeamerfont{title}\usebeamercolor[fg]{title}\inserttitle\par}
        \vskip0.3em
        {\usebeamerfont{subtitle}\usebeamercolor[fg]{subtitle}\insertsubtitle\par}
        \vskip1.2em
        \textcolor{ubbblue}{\rule{0.72\textwidth}{0.8pt}}
        \vskip1.2em
        {\small
            \textbf{\CourseName}\ (\CourseCode)\par
            \vskip0.4em
            \CourseProgram\par
            \CourseFaculty\par
            \CourseInstitution\par
            \vskip0.6em
            \CourseSemester\ \textperiodcentered\ \WeekLabel\par
            \vskip0.8em
            Dosen Pengampu: \CourseLecturer\par
        }
        \vskip1.0em
        \textcolor{ubbblue}{\rule{0.72\textwidth}{0.4pt}}
    \end{centering}
}

% =====================================================================
%  PERINTAH BANTU
% =====================================================================
\newcommand{\R}{\mathbb{R}}
\newcommand{\N}{\mathbb{N}}
\newcommand{\Z}{\mathbb{Z}}
\newcommand{\code}[1]{\texttt{#1}}

% Rujukan gambar dari buku: \bookfig[lebar]{NN}{namaberkas.png}
% Jalur relatif terhadap folder ppt/ (dikompilasi dari sana).
\newcommand{\bookfig}[3][0.80\linewidth]{%
    \includegraphics[width=#1]{../buku_ajar/chapters/#2/figures/#3}%
}

% Keterangan sumber gambar, sepadan dengan \sumber di buku
\newcommand{\figsumber}[1]{%
    \par\vspace{1pt}{\tiny\itshape Sumber: #1\par}%
}

% Penanda blok materi yang sering dipakai di slide
\newcommand{\slidejudul}[1]{{\large\bfseries\textcolor{ubbblue}{#1}}\par\vspace{2pt}}
\newcommand{\kunci}[1]{\textcolor{ubbred}{\textbf{#1}}}

%      \title{...} \subtitle{...} \setWeekInfo{Minggu n}{Topik}
%
%  Berkas minggu hanya berisi \setWeekInfo dan lingkungan frame.
%  Kompilasi dijalankan DARI DALAM folder ppt/ (lihat compile-ppt.bat),
%  karena gambar buku dirujuk dengan jalur relatif ../buku_ajar/.
% =====================================================================

% ---------------------------------------------------------------- bahasa
\usepackage[utf8]{inputenc}
\usepackage[T1]{fontenc}
\usepackage[indonesian]{babel}
\usepackage{lmodern}
\usepackage{microtype}

% ------------------------------------------------------------- matematika
\usepackage{amsmath,amssymb,mathtools}

% ----------------------------------------------------------------- tabel
\usepackage{booktabs}
\usepackage{array}
\usepackage{tabularx}

% --------------------------------------------------------------- grafik
\usepackage{tikz}
\usetikzlibrary{arrows.meta,positioning,shapes.geometric,calc,fit,backgrounds}

% ---------------------------------------------------------- kode program
\usepackage{listings}
\lstdefinestyle{slidecode}{
    basicstyle=\ttfamily\scriptsize,
    keywordstyle=\color{ubbblue}\bfseries,
    commentstyle=\color{ubbgreen}\itshape,
    stringstyle=\color{ubbred},
    numbers=left,
    numberstyle=\tiny\color{gray},
    numbersep=6pt,
    frame=single,
    rulecolor=\color{ubbblue!40},
    backgroundcolor=\color{ubbblue!3},
    breaklines=true,
    showstringspaces=false,
    tabsize=4,
    columns=fullflexible,
    keepspaces=true,
}
\lstset{style=slidecode}

% =====================================================================
%  WARNA INSTITUSI
% =====================================================================
\definecolor{ubbblue}{RGB}{12, 62, 120}      % warna utama
\definecolor{ubbteal}{RGB}{0, 118, 110}      % aksen kedua
\definecolor{ubbgreen}{RGB}{22, 122, 74}     % untuk blok contoh
\definecolor{ubbred}{RGB}{168, 42, 42}       % untuk blok peringatan
\definecolor{ubbgray}{RGB}{70, 76, 84}       % teks sekunder

% =====================================================================
%  IDENTITAS MATA KULIAH
% =====================================================================
\newcommand{\CourseName}{Pengantar Kriptografi}
\newcommand{\CourseCode}{MPFA014}
\newcommand{\CourseSemester}{Semester Ganjil}
\newcommand{\CourseProgram}{Program Studi Matematika}
\newcommand{\CourseFaculty}{Fakultas Matematika dan Ilmu Pengetahuan Alam}
\newcommand{\CourseInstitution}{Universitas Bale Bandung}
\newcommand{\CourseLecturer}{Cecep Suwanda, S.Si., M.Kom.}

% Label minggu, disetel tiap berkas lewat \setWeekInfo
\newcommand{\WeekLabel}{}
\newcommand{\WeekTopic}{}
\newcommand{\setWeekInfo}[2]{%
    \renewcommand{\WeekLabel}{#1}%
    \renewcommand{\WeekTopic}{#2}%
}

% =====================================================================
%  TEMA: header, footer, warna blok
% =====================================================================
\usetheme{default}
\setbeamertemplate{navigation symbols}{}
\setbeamertemplate{blocks}[rounded][shadow=false]

\setbeamercolor{palette primary}{bg=ubbblue, fg=white}
\setbeamercolor{palette secondary}{bg=ubbgray, fg=white}
\setbeamercolor{palette tertiary}{bg=ubbteal, fg=white}

\setbeamercolor{structure}{fg=ubbblue}
\setbeamercolor{normal text}{fg=black!92}
\setbeamercolor{alerted text}{fg=ubbred}

\setbeamercolor{block title}{bg=ubbblue, fg=white}
\setbeamercolor{block body}{bg=ubbblue!6, fg=black!92}
\setbeamercolor{block title alerted}{bg=ubbred, fg=white}
\setbeamercolor{block body alerted}{bg=ubbred!6, fg=black!92}
\setbeamercolor{block title example}{bg=ubbgreen, fg=white}
\setbeamercolor{block body example}{bg=ubbgreen!6, fg=black!92}

\setbeamercolor{itemize item}{fg=ubbblue}
\setbeamercolor{itemize subitem}{fg=ubbteal}
\setbeamercolor{enumerate item}{fg=ubbblue}

\setbeamerfont{headline}{size=\scriptsize}
\setbeamerfont{footline}{size=\scriptsize}
\setbeamerfont{block title}{size=\normalsize}

% Badan blok rata kiri (bukan rata kanan-kiri): teks slide lebih mudah dibaca
% dan baris terakhir yang pendek tidak lagi memicu "Underfull \hbox".
\addtobeamertemplate{block body begin}{\raggedright}{}
\addtobeamertemplate{block alerted body begin}{\raggedright}{}
\addtobeamertemplate{block example body begin}{\raggedright}{}

% Header: nama mata kuliah dan kode MK
\setbeamertemplate{headline}{%
    \begin{beamercolorbox}[wd=\paperwidth,ht=2.6ex,dp=1.2ex,leftskip=1em,rightskip=1em]{palette primary}%
        \usebeamerfont{headline}\CourseName\ \textperiodcentered\ \CourseCode%
        \hfill \CourseProgram%
    \end{beamercolorbox}%
}

% Footer: minggu dan nomor halaman
\setbeamertemplate{footline}{%
    \begin{beamercolorbox}[wd=\paperwidth,ht=2.4ex,dp=1.2ex,leftskip=1em,rightskip=1em]{palette secondary}%
        \usebeamerfont{footline}\WeekLabel\ \textperiodcentered\ \WeekTopic%
        \hfill \insertframenumber\,/\,\inserttotalframenumber%
    \end{beamercolorbox}%
}

% Halaman judul: judul, subjudul, lalu blok identitas mata kuliah
\setbeamertemplate{title page}{%
    \begin{centering}
        \vskip1.0em
        {\usebeamerfont{title}\usebeamercolor[fg]{title}\inserttitle\par}
        \vskip0.3em
        {\usebeamerfont{subtitle}\usebeamercolor[fg]{subtitle}\insertsubtitle\par}
        \vskip1.2em
        \textcolor{ubbblue}{\rule{0.72\textwidth}{0.8pt}}
        \vskip1.2em
        {\small
            \textbf{\CourseName}\ (\CourseCode)\par
            \vskip0.4em
            \CourseProgram\par
            \CourseFaculty\par
            \CourseInstitution\par
            \vskip0.6em
            \CourseSemester\ \textperiodcentered\ \WeekLabel\par
            \vskip0.8em
            Dosen Pengampu: \CourseLecturer\par
        }
        \vskip1.0em
        \textcolor{ubbblue}{\rule{0.72\textwidth}{0.4pt}}
    \end{centering}
}

% =====================================================================
%  PERINTAH BANTU
% =====================================================================
\newcommand{\R}{\mathbb{R}}
\newcommand{\N}{\mathbb{N}}
\newcommand{\Z}{\mathbb{Z}}
\newcommand{\code}[1]{\texttt{#1}}

% Rujukan gambar dari buku: \bookfig[lebar]{NN}{namaberkas.png}
% Jalur relatif terhadap folder ppt/ (dikompilasi dari sana).
\newcommand{\bookfig}[3][0.80\linewidth]{%
    \includegraphics[width=#1]{../buku_ajar/chapters/#2/figures/#3}%
}

% Keterangan sumber gambar, sepadan dengan \sumber di buku
\newcommand{\figsumber}[1]{%
    \par\vspace{1pt}{\tiny\itshape Sumber: #1\par}%
}

% Penanda blok materi yang sering dipakai di slide
\newcommand{\slidejudul}[1]{{\large\bfseries\textcolor{ubbblue}{#1}}\par\vspace{2pt}}
\newcommand{\kunci}[1]{\textcolor{ubbred}{\textbf{#1}}}

%      \title{...} \subtitle{...} \setWeekInfo{Minggu n}{Topik}
%
%  Berkas minggu hanya berisi \setWeekInfo dan lingkungan frame.
%  Kompilasi dijalankan DARI DALAM folder ppt/ (lihat compile-ppt.bat),
%  karena gambar buku dirujuk dengan jalur relatif ../buku_ajar/.
% =====================================================================

% ---------------------------------------------------------------- bahasa
\usepackage[utf8]{inputenc}
\usepackage[T1]{fontenc}
\usepackage[indonesian]{babel}
\usepackage{lmodern}
\usepackage{microtype}

% ------------------------------------------------------------- matematika
\usepackage{amsmath,amssymb,mathtools}

% ----------------------------------------------------------------- tabel
\usepackage{booktabs}
\usepackage{array}
\usepackage{tabularx}

% --------------------------------------------------------------- grafik
\usepackage{tikz}
\usetikzlibrary{arrows.meta,positioning,shapes.geometric,calc,fit,backgrounds}

% ---------------------------------------------------------- kode program
\usepackage{listings}
\lstdefinestyle{slidecode}{
    basicstyle=\ttfamily\scriptsize,
    keywordstyle=\color{ubbblue}\bfseries,
    commentstyle=\color{ubbgreen}\itshape,
    stringstyle=\color{ubbred},
    numbers=left,
    numberstyle=\tiny\color{gray},
    numbersep=6pt,
    frame=single,
    rulecolor=\color{ubbblue!40},
    backgroundcolor=\color{ubbblue!3},
    breaklines=true,
    showstringspaces=false,
    tabsize=4,
    columns=fullflexible,
    keepspaces=true,
}
\lstset{style=slidecode}

% =====================================================================
%  WARNA INSTITUSI
% =====================================================================
\definecolor{ubbblue}{RGB}{12, 62, 120}      % warna utama
\definecolor{ubbteal}{RGB}{0, 118, 110}      % aksen kedua
\definecolor{ubbgreen}{RGB}{22, 122, 74}     % untuk blok contoh
\definecolor{ubbred}{RGB}{168, 42, 42}       % untuk blok peringatan
\definecolor{ubbgray}{RGB}{70, 76, 84}       % teks sekunder

% =====================================================================
%  IDENTITAS MATA KULIAH
% =====================================================================
\newcommand{\CourseName}{Pengantar Kriptografi}
\newcommand{\CourseCode}{MPFA014}
\newcommand{\CourseSemester}{Semester Ganjil}
\newcommand{\CourseProgram}{Program Studi Matematika}
\newcommand{\CourseFaculty}{Fakultas Matematika dan Ilmu Pengetahuan Alam}
\newcommand{\CourseInstitution}{Universitas Bale Bandung}
\newcommand{\CourseLecturer}{Cecep Suwanda, S.Si., M.Kom.}

% Label minggu, disetel tiap berkas lewat \setWeekInfo
\newcommand{\WeekLabel}{}
\newcommand{\WeekTopic}{}
\newcommand{\setWeekInfo}[2]{%
    \renewcommand{\WeekLabel}{#1}%
    \renewcommand{\WeekTopic}{#2}%
}

% =====================================================================
%  TEMA: header, footer, warna blok
% =====================================================================
\usetheme{default}
\setbeamertemplate{navigation symbols}{}
\setbeamertemplate{blocks}[rounded][shadow=false]

\setbeamercolor{palette primary}{bg=ubbblue, fg=white}
\setbeamercolor{palette secondary}{bg=ubbgray, fg=white}
\setbeamercolor{palette tertiary}{bg=ubbteal, fg=white}

\setbeamercolor{structure}{fg=ubbblue}
\setbeamercolor{normal text}{fg=black!92}
\setbeamercolor{alerted text}{fg=ubbred}

\setbeamercolor{block title}{bg=ubbblue, fg=white}
\setbeamercolor{block body}{bg=ubbblue!6, fg=black!92}
\setbeamercolor{block title alerted}{bg=ubbred, fg=white}
\setbeamercolor{block body alerted}{bg=ubbred!6, fg=black!92}
\setbeamercolor{block title example}{bg=ubbgreen, fg=white}
\setbeamercolor{block body example}{bg=ubbgreen!6, fg=black!92}

\setbeamercolor{itemize item}{fg=ubbblue}
\setbeamercolor{itemize subitem}{fg=ubbteal}
\setbeamercolor{enumerate item}{fg=ubbblue}

\setbeamerfont{headline}{size=\scriptsize}
\setbeamerfont{footline}{size=\scriptsize}
\setbeamerfont{block title}{size=\normalsize}

% Badan blok rata kiri (bukan rata kanan-kiri): teks slide lebih mudah dibaca
% dan baris terakhir yang pendek tidak lagi memicu "Underfull \hbox".
\addtobeamertemplate{block body begin}{\raggedright}{}
\addtobeamertemplate{block alerted body begin}{\raggedright}{}
\addtobeamertemplate{block example body begin}{\raggedright}{}

% Header: nama mata kuliah dan kode MK
\setbeamertemplate{headline}{%
    \begin{beamercolorbox}[wd=\paperwidth,ht=2.6ex,dp=1.2ex,leftskip=1em,rightskip=1em]{palette primary}%
        \usebeamerfont{headline}\CourseName\ \textperiodcentered\ \CourseCode%
        \hfill \CourseProgram%
    \end{beamercolorbox}%
}

% Footer: minggu dan nomor halaman
\setbeamertemplate{footline}{%
    \begin{beamercolorbox}[wd=\paperwidth,ht=2.4ex,dp=1.2ex,leftskip=1em,rightskip=1em]{palette secondary}%
        \usebeamerfont{footline}\WeekLabel\ \textperiodcentered\ \WeekTopic%
        \hfill \insertframenumber\,/\,\inserttotalframenumber%
    \end{beamercolorbox}%
}

% Halaman judul: judul, subjudul, lalu blok identitas mata kuliah
\setbeamertemplate{title page}{%
    \begin{centering}
        \vskip1.0em
        {\usebeamerfont{title}\usebeamercolor[fg]{title}\inserttitle\par}
        \vskip0.3em
        {\usebeamerfont{subtitle}\usebeamercolor[fg]{subtitle}\insertsubtitle\par}
        \vskip1.2em
        \textcolor{ubbblue}{\rule{0.72\textwidth}{0.8pt}}
        \vskip1.2em
        {\small
            \textbf{\CourseName}\ (\CourseCode)\par
            \vskip0.4em
            \CourseProgram\par
            \CourseFaculty\par
            \CourseInstitution\par
            \vskip0.6em
            \CourseSemester\ \textperiodcentered\ \WeekLabel\par
            \vskip0.8em
            Dosen Pengampu: \CourseLecturer\par
        }
        \vskip1.0em
        \textcolor{ubbblue}{\rule{0.72\textwidth}{0.4pt}}
    \end{centering}
}

% =====================================================================
%  PERINTAH BANTU
% =====================================================================
\newcommand{\R}{\mathbb{R}}
\newcommand{\N}{\mathbb{N}}
\newcommand{\Z}{\mathbb{Z}}
\newcommand{\code}[1]{\texttt{#1}}

% Rujukan gambar dari buku: \bookfig[lebar]{NN}{namaberkas.png}
% Jalur relatif terhadap folder ppt/ (dikompilasi dari sana).
\newcommand{\bookfig}[3][0.80\linewidth]{%
    \includegraphics[width=#1]{../buku_ajar/chapters/#2/figures/#3}%
}

% Keterangan sumber gambar, sepadan dengan \sumber di buku
\newcommand{\figsumber}[1]{%
    \par\vspace{1pt}{\tiny\itshape Sumber: #1\par}%
}

% Penanda blok materi yang sering dipakai di slide
\newcommand{\slidejudul}[1]{{\large\bfseries\textcolor{ubbblue}{#1}}\par\vspace{2pt}}
\newcommand{\kunci}[1]{\textcolor{ubbred}{\textbf{#1}}}


\title{Landasan Matematika untuk Kriptografi}
\subtitle{Aritmetika modulo, invers modulo, bilangan prima, dan entropi}
\setWeekInfo{Minggu 3}{Landasan matematika}

\begin{document}

% ----------------------------------------------------------- halaman judul
\begin{frame}[plain]
    \titlepage
\end{frame}

% ------------------------------------------------------------------ peta
\begin{frame}{Peta Materi Pertemuan Ini}
    \begin{columns}[T,onlytextwidth]
        \begin{column}{0.54\textwidth}
            \slidejudul{Enam perkakas}
            \begin{enumerate}
                \item Aritmetika modulo dan kekongruenan
                \item Relatif prima dan PBB
                \item Algoritma Euclidean
                \item Invers modulo
                \item Bilangan prima
                \item Entropi dan teori informasi
            \end{enumerate}
        \end{column}
        \begin{column}{0.42\textwidth}
            \begin{block}{Referensi}
                \footnotesize
                Bab 3 buku ajar; Munir, \textit{IF4020 Kriptografi},
                ITB, 2026 (modul 02)
            \end{block}
        \end{column}
    \end{columns}
\end{frame}

% ---------------------------------------------------- capaian pembelajaran
\begin{frame}{Capaian Pembelajaran}
    \begin{block}{Sub-CPMK yang dibebankan}
        \footnotesize
        \textbf{Sub-CPMK 2.3} --- Menjelaskan landasan matematika dasar untuk
        kriptografi, yaitu aritmetika modulo, relatif prima, invers modulo,
        bilangan prima, serta entropi dan teori informasi.
    \end{block}
    \vspace{0.2cm}
    \begin{exampleblock}{Setelah pertemuan ini, mahasiswa mampu}
        \begin{itemize}
            \item menjelaskan aturan sisa bagi dan kekongruenan;
            \item menjelaskan syarat dua bilangan relatif prima;
            \item mengidentifikasi invers modulo dan memeriksa hasilnya;
            \item menjelaskan peran bilangan prima pada algoritma kunci publik;
            \item menjelaskan entropi dan menafsirkan nilainya.
        \end{itemize}
    \end{exampleblock}
\end{frame}

% =====================================================================
\section{Aritmetika Modulo dan Kekongruenan}
% =====================================================================

\begin{frame}{Mengapa Kriptografi Menuntut Matematika}
    Caesar cipher cukup dijelaskan dengan pergeseran huruf; algoritma kunci
    publik tidak.
    \vspace{0.1cm}
    \footnotesize
    \begin{tabularx}{\linewidth}{@{}>{\bfseries\raggedright\arraybackslash}p{4.8cm}
                                  >{\raggedright\arraybackslash}X@{}}
        \toprule
        Perkakas matematika & Dipakai untuk \\
        \midrule
        Aritmetika modulo dan kekongruenan & Cipher substitusi seperti Caesar,
        Vigen\`ere, dan affine. \\
        \addlinespace
        PBB, relatif prima, invers modulo & Pemilihan kunci dan dekripsi RSA. \\
        \addlinespace
        Bilangan prima & Pembangkitan kunci pada algoritma kunci publik. \\
        \addlinespace
        Entropi dan teori informasi & Mengukur kebocoran informasi dan menilai
        ketahanan cipherteks. \\
        \bottomrule
    \end{tabularx}
    \vspace{0.15cm}
    \begin{alertblock}{Benang merah}
        Tanpa landasan ini, algoritma kriptografi hanya tampak sebagai langkah
        yang dihafal.
    \end{alertblock}
\end{frame}

\begin{frame}{Aritmetika Modulo: Sisa Bagi yang Selalu Nonnegatif}
    Jika $a$ bilangan bulat dan $n$ bilangan bulat positif, maka $a \bmod n$
    adalah sisa pembagian $a$ oleh $n$.
    \[ a = qn + r, \qquad 0 \le r < n. \]
    Nilai $r$ inilah yang disebut $a \bmod n$, dan syarat $0 \le r < n$ membuat
    hasilnya \kunci{tunggal}.
    \vspace{0.1cm}
    \begin{columns}[T,onlytextwidth]
        \begin{column}{0.48\textwidth}
            \begin{block}{Contoh}
                \begin{itemize}
                    \item $23 \bmod 5 = 3$, karena $23 = 4 \cdot 5 + 3$.
                    \item $-41 \bmod 9 = 4$, karena
                          $-41 = (-5)\cdot 9 + 4$.
                \end{itemize}
            \end{block}
        \end{column}
        \begin{column}{0.48\textwidth}
            \begin{exampleblock}{Membaca hasilnya}
                Pada contoh kedua, $-45$ adalah kelipatan 9 terbesar yang tidak
                melebihi $-41$; selisihnya $4$ menjadi sisa.
            \end{exampleblock}
        \end{column}
    \end{columns}
\end{frame}

\begin{frame}{Contoh Modulo Bilangan Negatif}
    \begin{columns}[T,onlytextwidth]
        \begin{column}{0.52\textwidth}
            Aturannya tidak berubah: jawaban harus berada pada $0 \le r < n$.
            Karena itu hasil modulo bilangan negatif \kunci{tidak boleh negatif}.
            \vspace{0.2cm}
            \begin{itemize}
                \item $-41 \bmod 9 = 4$, sebab
                      $-41 = (-5)\cdot 9 + 4$.
                \item $-24 \bmod 11 = 9$, sebab
                      $-24 = (-3)\cdot 11 + 9$.
            \end{itemize}
        \end{column}
        \begin{column}{0.44\textwidth}
            \begin{exampleblock}{Pemeriksaan}
                Pada contoh kedua, $(-3)\cdot 11 = -33$ dan
                $-24 - (-33) = 9$. Mengalikan kembali seperti ini adalah
                kebiasaan memeriksa yang baik.
            \end{exampleblock}
        \end{column}
    \end{columns}
\end{frame}

\begin{frame}{Kekongruenan: Pernyataan tentang Selisih}
    Dua bilangan $a$ dan $b$ disebut \textbf{kongruen modulo $n$}, ditulis
    $a \equiv b \pmod n$, jika $n$ membagi habis $a - b$.
    \vspace{0.15cm}
    \begin{columns}[T,onlytextwidth]
        \begin{column}{0.48\textwidth}
            \begin{itemize}
                \item $17 \equiv 2 \pmod 3$, karena $3 \mid (17 - 2)$.
                \item $-7 \equiv 15 \pmod{11}$, karena
                      $11 \mid (-7 - 15)$.
            \end{itemize}
        \end{column}
        \begin{column}{0.48\textwidth}
            \begin{alertblock}{Dua notasi yang berbeda}
                $a \bmod n$ adalah sebuah \emph{bilangan}, yaitu sisanya;
                $a \equiv b \pmod n$ adalah \emph{pernyataan} benar atau salah.
            \end{alertblock}
        \end{column}
    \end{columns}
\end{frame}

\begin{frame}{Sifat yang Mengikuti Definisi Kekongruenan}
    Bila $a \equiv b \pmod n$, maka menambahkan atau mengalikan bilangan yang
    sama pada kedua ruas tidak mengubah kebenarannya.
    \vspace{0.2cm}
    \begin{block}{Tiga akibat langsung}
        \begin{itemize}
            \item $a + c \equiv b + c \pmod n$ untuk setiap bilangan bulat $c$;
            \item $a \cdot c \equiv b \cdot c \pmod n$ untuk setiap bilangan
                  bulat $c$;
            \item $a^k \equiv b^k \pmod n$ untuk setiap bilangan bulat positif
                  $k$.
        \end{itemize}
    \end{block}
    \vspace{0.15cm}
    \begin{exampleblock}{Contoh pemakaian}
        Karena $17 \equiv 2 \pmod 3$, maka $17 \cdot 5 \equiv 2 \cdot 5
        \pmod 3$ dan $17^2 \equiv 2^2 \pmod 3$.
    \end{exampleblock}
\end{frame}

% =====================================================================
\section{Relatif Prima dan Invers Modulo}
% =====================================================================

\begin{frame}{Relatif Prima dan Pembagi Persekutuan Terbesar}
    \begin{columns}[T,onlytextwidth]
        \begin{column}{0.50\textwidth}
            \begin{block}{Definisi}
                PBB dari $a$ dan $b$ adalah bilangan bulat positif terbesar
                yang membagi habis keduanya. Bila $\text{PBB}(a,b) = 1$, maka
                $a$ dan $b$ disebut \textbf{relatif prima}.
            \end{block}
            \begin{itemize}
                \item $23$ dan $40$ relatif prima karena
                      $\text{PBB}(23,40) = 1$.
                \item $48$ dan $18$ tidak relatif prima karena
                      $\text{PBB}(48,18) = 6$.
            \end{itemize}
        \end{column}
        \begin{column}{0.46\textwidth}
            \begin{alertblock}{Kaitannya dengan invers}
                Relatif prima adalah syarat keberadaan invers modulo: invers
                $a$ modulo $n$ hanya ada bila $\text{PBB}(a,n) = 1$.
            \end{alertblock}
        \end{column}
    \end{columns}
\end{frame}

\begin{frame}{Algoritma Euclidean}
    PBB dihitung dengan pembagian berulang sampai sisanya nol. Sisa bukan nol
    terakhir adalah PBB-nya.
    \begin{columns}[T,onlytextwidth]
        \begin{column}{0.48\textwidth}
            \[
            \begin{aligned}
                48 &= 2 \cdot 18 + 12\\
                18 &= 1 \cdot 12 + 6\\
                12 &= 2 \cdot 6 + 0
            \end{aligned}
            \]
            Sisa bukan nol terakhir adalah $6$, jadi
            \kunci{$\text{PBB}(48,18) = 6$}.
        \end{column}
        \begin{column}{0.48\textwidth}
            \begin{block}{Cara kerja}
                Setiap langkah mengganti pasangan bilangan dengan pasangan
                baru yang lebih kecil, yaitu pembagi dan sisanya. Proses
                berhenti ketika sisa menjadi nol.
            \end{block}
        \end{column}
    \end{columns}
\end{frame}

\begin{frame}{Invers Modulo}
    Invers modulo $a$ modulo $n$ adalah bilangan $x$ yang memenuhi
    \[ a \cdot x \equiv 1 \pmod n. \]
    Invers itu ada jika dan hanya jika $a$ dan $n$ relatif prima.
    \vspace{0.1cm}
    \begin{columns}[T,onlytextwidth]
        \begin{column}{0.50\textwidth}
            \begin{exampleblock}{Dua contoh dari buku ajar}
                \begin{itemize}
                    \item $4^{-1} \bmod 9 \equiv 7 \pmod 9$, karena
                          $4 \cdot 7 = 28 = 3 \cdot 9 + 1$.
                    \item $23^{-1} \bmod 10 \equiv -3 \equiv 7 \pmod{10}$,
                          karena $23 \cdot (-3) = -69 = (-7)\cdot 10 + 1$.
                \end{itemize}
            \end{exampleblock}
        \end{column}
        \begin{column}{0.46\textwidth}
            \begin{alertblock}{Bukan kebalikan biasa}
                Pada modulo 9, $4^{-1}$ bukan $1/4$ melainkan $7$. Invers
                selalu diperiksa dengan mengalikan kembali lalu mengurangi
                modulo $n$.
            \end{alertblock}
        \end{column}
    \end{columns}
\end{frame}

\begin{frame}{Contoh Terhitung --- Invers $23^{-1} \bmod 10$ (1/2)}
    \footnotesize
    Syarat diperiksa lebih dahulu dengan Algoritma Euclidean. Setiap baris
    mencatat operasi pembagian dan sisanya.
    \vspace{0.15cm}
    \begin{tabularx}{\linewidth}{@{}>{\bfseries}p{1.0cm}
                                  >{\raggedright\arraybackslash}p{7.0cm}
                                  >{\raggedright\arraybackslash}X@{}}
        \toprule
        No. & Operasi & Hasil \\
        \midrule
        1 & $23 = 2 \cdot 10 + 3$ & sisa $3$ \\
        \addlinespace
        2 & $10 = 3 \cdot 3 + 1$ & sisa $1$ \\
        \addlinespace
        3 & $3 = 3 \cdot 1 + 0$ & sisa $0$ \\
        \bottomrule
    \end{tabularx}
    \vspace{0.25cm}
    \begin{block}{Hasil pemeriksaan syarat}
        Sisa bukan nol terakhir adalah $1$, sehingga
        $\text{PBB}(23,10) = 1$. Karena kedua bilangan relatif prima, invers
        $23^{-1} \bmod 10$ ada.
    \end{block}
\end{frame}

\begin{frame}{Contoh Terhitung --- Invers $23^{-1} \bmod 10$ (2/2)}
    \footnotesize
    Koefisien dicari dengan substitusi balik dari sisa terakhir ke atas.
    \vspace{0.15cm}
    \begin{tabularx}{\linewidth}{@{}>{\bfseries}p{1.0cm}
                                  >{\raggedright\arraybackslash}p{8.2cm}
                                  >{\raggedright\arraybackslash}X@{}}
        \toprule
        No. & Operasi & Hasil \\
        \midrule
        4 & $1 = 10 - 3 \cdot 3$ & bentuk awal \\
        \addlinespace
        5 & $3 = 23 - 2 \cdot 10$ & siap disubstitusikan \\
        \addlinespace
        6 & $1 = 10 - 3(23 - 2 \cdot 10) = 7 \cdot 10 - 3 \cdot 23$
          & koefisien $23$ adalah $-3$ \\
        \addlinespace
        7 & $x \equiv -3 \pmod{10}$ lalu dinormalkan & $x = 7$ \\
        \bottomrule
    \end{tabularx}
    \vspace{0.2cm}
    \begin{exampleblock}{Pemeriksaan hasil}
        $23 \cdot 7 = 161 = 16 \cdot 10 + 1$, jadi
        $23 \cdot 7 \equiv 1 \pmod{10}$. Hasil ini setara dengan
        $23 \cdot (-3) = -69 = (-7)\cdot 10 + 1$ pada buku ajar.
    \end{exampleblock}
\end{frame}

% =====================================================================
\section{Bilangan Prima}
% =====================================================================

\begin{frame}{Bilangan Prima}
    \begin{columns}[T,onlytextwidth]
        \begin{column}{0.50\textwidth}
            \begin{block}{Definisi}
                Bilangan prima adalah bilangan bulat lebih besar dari 1 yang
                hanya habis dibagi oleh 1 dan dirinya sendiri.
            \end{block}
            \begin{itemize}
                \item $2, 3, 5, 7, 11, 13$ adalah bilangan prima.
                \item $9$ bukan prima karena $9 = 3 \cdot 3$.
                \item $1$ bukan prima karena tidak lebih besar dari 1.
            \end{itemize}
        \end{column}
        \begin{column}{0.46\textwidth}
            \begin{alertblock}{Perannya pada kriptografi}
                Bilangan prima dipakai pada algoritma kunci publik. Sifat
                \"hanya habis dibagi 1 dan dirinya sendiri\" menjadi bahan
                pembangkitan kunci.
            \end{alertblock}
        \end{column}
    \end{columns}
\end{frame}

% =====================================================================
\section{Teori Informasi dan Entropi}
% =====================================================================

\begin{frame}{Teori Informasi: Ukuran dalam Bit}
    Teori informasi mempelajari kuantisasi, penyimpanan, transmisi, dan
    pengolahan informasi. Bagi kriptografi, ia memberi ukuran seberapa besar
    informasi yang dibocorkan sebuah pesan.
    \vspace{0.15cm}
    \begin{columns}[T,onlytextwidth]
        \begin{column}{0.50\textwidth}
            \begin{block}{Banyaknya bit bergantung pada pilihan}
                \begin{itemize}
                    \item $1$ bit cukup untuk jenis kelamin, karena ada dua
                          kemungkinan.
                    \item $3$ bit untuk nama hari, karena $2^3 = 8 \ge 7$.
                    \item $4$ bit untuk angka $0$ sampai $9$, karena
                          $2^4 = 16 \ge 10$.
                \end{itemize}
            \end{block}
        \end{column}
        \begin{column}{0.46\textwidth}
            \begin{exampleblock}{Penyimpanan}
                Pemampatan data mengurangi ukuran ruang penyimpanan dengan
                membuang bagian pesan yang dapat diduga.
            \end{exampleblock}
        \end{column}
    \end{columns}
\end{frame}

\begin{frame}{Entropi Shannon}
    \begin{block}{Rumus entropi}
        \[ H(X) = -\sum_{i=1}^{n} p(x_i) \log_2 p(x_i), \]
        dengan $X$ variabel acak pesan, $x_i$ simbol ke-$i$, $n$ banyaknya
        simbol, dan $p(x_i)$ peluang kemunculan simbol $x_i$.
    \end{block}
    \vspace{0.15cm}
    Entropi mengukur ketidakpastian, yaitu jumlah rata-rata informasi di
    dalam pesan, dan dinyatakan dalam bit.
    \vspace{0.1cm}
    \begin{itemize}
        \item Semakin acak suatu data, semakin tinggi entropinya.
        \item Cipherteks adalah pesan dengan entropi tinggi.
    \end{itemize}
\end{frame}

\begin{frame}{Contoh Entropi Sebuah Pesan Pendek}
    \footnotesize
    Pesan $X = \code{AABBCBDB}$ memuat 8 simbol dengan $p(A) = 2/8$,
    $p(B) = 4/8$, serta $p(C) = p(D) = 1/8$. Substitusi ke rumus memberi
    \[
    \begin{aligned}
        H(X) &= -\left(\tfrac14 \log_2 \tfrac14 + \tfrac12 \log_2 \tfrac12
                 + \tfrac18 \log_2 \tfrac18 + \tfrac18 \log_2 \tfrac18\right)\\
             &= -\left(\tfrac14(-2) + \tfrac12(-1) + \tfrac18(-3)
                 + \tfrac18(-3)\right)\\
             &= 1{,}75\ \text{bit/simbol}.
    \end{aligned}
    \]
    \begin{exampleblock}{Tafsir}
        Entropi maksimum untuk empat simbol adalah $\log_2 4 = 2$ bit.
        Sebaran yang tidak seragam ini menghasilkan $1{,}75$ bit/simbol,
        yaitu lebih rendah daripada maksimumnya.
    \end{exampleblock}
\end{frame}

\begin{frame}{Rentang Nilai Entropi}
    \begin{columns}[T,onlytextwidth]
        \begin{column}{0.54\textwidth}
            \[ 0 \le H(X) \le \log_2 n. \]
            \begin{itemize}
                \item \textbf{Minimum} $H = 0$ bit terjadi bila hanya ada satu
                      simbol dengan peluang 1: $H = -1 \cdot \log_2 1 = 0$.
                \item \textbf{Maksimum} $H = \log_2 n$ bit terjadi bila semua
                      simbol berpeluang sama $1/n$. Untuk empat simbol dengan
                      peluang $0{,}25$, diperoleh $2$ bit.
            \end{itemize}
        \end{column}
        \begin{column}{0.42\textwidth}
            \centering
            \begin{tikzpicture}[font=\scriptsize]
                \draw[-{Latex[length=2.2mm]}, thick] (0,0) -- (5.2,0);
                \draw[ubbblue, line width=5pt] (0.25,0.22) -- (4.95,0.22);
                \node[below] at (0.25,-0.02) {$0$};
                \node[below] at (4.95,-0.02) {$\log_2 n$};
                \node[above, align=center] at (1.15,0.50)
                    {satu simbol};
                \node[above, align=center] at (4.15,0.50)
                    {semua peluang sama};
            \end{tikzpicture}
            \figsumber{Rentang entropi, diringkas dari buku ajar Bab 3.}
        \end{column}
    \end{columns}
\end{frame}

\begin{frame}{Entropi Maksimum Alfabet Tertentu}
    Bila seluruh 26 huruf A--Z berpeluang sama, yaitu $p_i = 1/26$, maka
    \[
        H_{\max} = -\sum_{i=1}^{26} \frac{1}{26} \log_2 \frac{1}{26}
                 = \log_2 26 = 4{,}7004\ \text{bit/karakter}.
    \]
    \begin{columns}[T,onlytextwidth]
        \begin{column}{0.48\textwidth}
            \begin{exampleblock}{Pembanding}
                Untuk 256 karakter ASCII, entropi maksimumnya adalah
                $\log_2 256 = 8$ bit/karakter.
            \end{exampleblock}
        \end{column}
        \begin{column}{0.48\textwidth}
            \begin{alertblock}{Teks nyata tidak seragam}
                Pesan bermakna tidak memakai huruf dengan peluang sama,
                sehingga entropinya berada di bawah $4{,}7004$ bit/karakter.
            \end{alertblock}
        \end{column}
    \end{columns}
\end{frame}

\begin{frame}{Studi Kasus: Entropi Berita Kota Ternate}
    \footnotesize
    \begin{columns}[T,onlytextwidth]
        \begin{column}{0.52\textwidth}
            Kutipan berita memuat 296 huruf setelah spasi dan tanda baca
            dibuang, dengan 21 huruf berbeda. Sebarannya tidak rata: huruf A
            muncul 59 kali, N 34 kali, dan E 29 kali.
            \vspace{0.1cm}
            \[ H(P) = 3{,}8988\ \text{bit/karakter} \approx 82{,}95\%. \]
            Setelah dienkripsi, entropi cipherteks menjadi
            \[ H(C) = 4{,}3035\ \text{bit/karakter} \approx 91{,}55\%. \]
            Kedua persentase dihitung terhadap $\log_2 26 = 4{,}7004$.
        \end{column}
        \begin{column}{0.44\textwidth}
            \centering
            \bookfig[1.0\linewidth]{03}{frekuensi-huruf-ternate.png}
            \figsumber{Munir, \textit{IF4020 Kriptografi}, ITB.}
        \end{column}
    \end{columns}
\end{frame}

\begin{frame}{Entropi dan Keamanan Cipherteks}
    \begin{block}{Makin tinggi entropi, makin sulit dipecahkan}
        Cipherteks yang sebarannya masih menyerupai plainteks menyimpan pola
        bahasa dan rentan terhadap analisis frekuensi.
    \end{block}
    \vspace{0.2cm}
    \begin{columns}[T,onlytextwidth]
        \begin{column}{0.48\textwidth}
            \begin{exampleblock}{Cipherteks yang baik}
                Sebaran hurufnya lebih rata, sehingga pola plainteks tidak
                lagi tampak.
            \end{exampleblock}
        \end{column}
        \begin{column}{0.48\textwidth}
            \begin{alertblock}{Pelajaran dari studi kasus}
                Entropi naik dari $3{,}8988$ menjadi $4{,}3035$
                bit/karakter setelah enkripsi meratakan sebaran huruf.
            \end{alertblock}
        \end{column}
    \end{columns}
\end{frame}

% =====================================================================
\section{Rangkuman dan Latihan}
% =====================================================================

\begin{frame}{Rangkuman (1/2)}
    \footnotesize
    \begin{itemize}
        \item $a \bmod n$ adalah sisa pembagian dengan $0 \le r < n$;
              untuk bilangan negatif hasilnya tetap nonnegatif, misalnya
              $-41 \bmod 9 = 4$ dan $-24 \bmod 11 = 9$.
        \item $a \equiv b \pmod n$ berarti $n$ membagi habis $a - b$;
              notasi ini menyatakan pernyataan, bukan bilangan.
        \item Penjumlahan, perkalian, dan pemangkatan pada kedua ruas
              mempertahankan kebenaran kekongruenan.
        \item Dua bilangan relatif prima bila PBB keduanya 1, misalnya
              $\text{PBB}(23,40) = 1$.
        \item Algoritma Euclidean mencari PBB melalui pembagian berulang;
              contohnya $\text{PBB}(48,18) = 6$.
        \item Invers modulo $x = a^{-1} \bmod n$ memenuhi
              $a \cdot x \equiv 1 \pmod n$ dan hanya ada bila
              $\text{PBB}(a,n) = 1$; contohnya $4^{-1} \bmod 9 \equiv 7$.
    \end{itemize}
\end{frame}

\begin{frame}{Rangkuman (2/2)}
    \footnotesize
    \begin{itemize}
        \item Bilangan prima adalah bilangan lebih besar dari 1 yang hanya
              habis dibagi 1 dan dirinya sendiri, misalnya
              $2, 3, 5, 7, 11, 13$.
        \item Bilangan prima dipakai pada algoritma kunci publik.
        \item Entropi Shannon
              $H(X) = -\sum_i p(x_i)\log_2 p(x_i)$ mengukur ketidakpastian
              pesan dalam bit.
        \item Rentang entropi adalah $0 \le H(X) \le \log_2 n$; nilai minimum
              terjadi pada satu simbol berpeluang 1, nilai maksimum pada
              sebaran seragam.
        \item Entropi maksimum alfabet 26 huruf adalah
              $\log_2 26 = 4{,}7004$ bit/karakter, sedangkan ASCII memberi
              $8$ bit/karakter.
        \item Pada berita Kota Ternate, entropi naik dari $3{,}8988$ menjadi
              $4{,}3035$ bit/karakter setelah enkripsi.
    \end{itemize}
\end{frame}

\begin{frame}{Latihan}
    \footnotesize
    \begin{enumerate}
        \item Hitunglah nilai dari $-24 \bmod 11$.
        \item Tentukan invers modulo dari $12 \bmod 5$, lalu periksa hasilnya
              dengan mengalikan kembali.
        \item Sebuah pesan memiliki huruf A, B, C, dan D dengan peluang
              masing-masing $0{,}25$. Hitunglah entropinya.
        \item Tentukan $\text{PBB}(48,18)$ dengan Algoritma Euclidean.
        \item Jelaskan mengapa invers modulo $a^{-1} \bmod n$ hanya ada jika
              $\text{PBB}(a,n) = 1$.
        \item Sebuah pesan 200 huruf memiliki entropi $4{,}2$ bit/karakter.
              Berapa persen entropi maksimum alfabet 26 huruf yang dicapainya,
              dan apa artinya bagi ketahanan cipherteks terhadap analisis
              frekuensi?
    \end{enumerate}
\end{frame}

\begin{frame}{Yang Harus Diingat}
    \begin{block}{Empat butir kunci}
        \begin{itemize}
            \item Sisa bagi selalu memenuhi $0 \le r < n$, termasuk ketika
                  bilangan yang dibagi negatif.
            \item Invers modulo hanya ada untuk pasangan yang relatif prima,
                  dan hasilnya diperiksa dengan perkalian.
            \item Bilangan prima menjadi bahan algoritma kunci publik.
            \item Entropi mengukur ketidakpastian; enkripsi yang meratakan
                  sebaran huruf menaikkan entropi cipherteks.
        \end{itemize}
    \end{block}
    \vspace{0.3cm}
    \begin{exampleblock}{Pertemuan berikutnya}
        \textbf{Minggu 4 --- Kriptografi klasik I}
    \end{exampleblock}
\end{frame}

\end{document}
